% !TeX program=lualatex

% Based on a "Simple LaTex Conference Poster Template" by @klb2:
% https://github.com/klb2/scientific-poster-template
%
\documentclass[25pt, a0paper, portrait]{littlefs-poster}
%\documentclass[25pt, a0paper, portrait]{tikzposter}
\geometry{paperheight=48in,paperwidth=36in} % for different sizes than A paper

% typesetting things
\usepackage[english]{babel}
\usepackage{microtype}
\usepackage[T1]{fontenc}
\usepackage[utf8]{inputenc}
\usepackage{pslatex}

% use sans-serif font
%\usepackage{lmodern}
\usepackage{amsmath,amsthm}
%\usepackage{newtxsf}
%\usepackage[italic]{mathastext}
%\usepackage{isomath}
\renewcommand{\familydefault}{\sfdefault}

% misc packages

%\usepackage[defaultsans,default]{opensans}
%\usepackage{fontspec} % use with lualatex or xelatex
%\usepackage{arev} % math font
%\setmainfont{Open Sans}

\usepackage{blindtext}

\usepackage{booktabs}
\usepackage{multirow}

\usetikzlibrary{positioning}

% old habits die hard
\def\lt{<}
\def\gt{>}





%%%% poster %%%
\title{\bf \fontsize{72}{80}\selectfont
    littlefs: rbyd B-trees: \\
    Incremental B-trees when RAM $\ll$ block size
}
\author{\LARGE
    Christopher Haster % \quad\textit{littlefs}
}
%\author{First Author\affilmark{1}, Second Author\affilmark{1,2}, and
%    Third Author\affilmark{2}}
%\institute{%
%\affilmark{1} Department of Science, University of City, Country\\
%\affilmark{2} Department of Engineering, Other University, Different Country
%}


%\titlegraphic{\includegraphics[height=.08\paperheight]{logo.png}}

%\usetitlestyle[titlebackgroundheight=.078\paperheight,titlegraphictotitledistance=1cm]{conferenceTitleStyle}


\begin{document}
\maketitle
\node[anchor=east] at ($(TP@title.east) +(-56pt,0pt)$) {
    \def\littlefsicoscale{7pt}
    \def\littlefsicopadding{2}
    \input{littlefs-ico}
};


% Problem statement
\begin{columns}
\column{0.5}
\block{Problem}{\Large
    Microcontrollers are very small, flash is very dense.

    How do you build a B-tree if you can't fit a block in RAM?

    \begin{enumerate}
    \item Traditional B-trees read a full block into RAM.
    \item Log-backed B-trees usually reconstruct logs in RAM.
    \end{enumerate}
}
\column{0.5}
\block[bodyinnersep=0pt]{}{\Large
    \begin{tikzfigure}
    \centering
    % MCU memory table
    \begin{tabular}{l@{\hspace{1em}}r}
        \toprule
            & MCU memory \\
        \midrule
        ROM & {$\sim$}256KiB--1MiB \\
        RAM & {$\sim$}32KiB--128KiB \\
        \bottomrule
    \end{tabular}
    \quad
    % Flash memory table
    \begin{tabular}{l@{\hspace{1em}}r@{\hspace{1em}}r}
        \toprule
            & NOR flash & NAND flash \\
        \midrule
        disk size & {$\sim$}128KiB--256MiB & {$\sim$}64MiB--2GiB \\
        block size & {$\sim$}4KiB--8KiB & {$\sim$}128KiB--1MiB \\
        \bottomrule
    \end{tabular}
    \end{tikzfigure}
}
\end{columns}


% Dhara and red-black-yellow Dhara trees
\begin{columns}
\column{0.5}
\block{Dhara trees \vphantom{y}}{\Large
}
\column{0.5}
\block{Red-black-yellow Dhara (rbyd) trees}{\Large
}
\end{columns}



% Main rbyd B-tree graphic!

% some defines
\def\bnodewidth{(1/(2*(1.5*7-0.5)))*\linewidth}
\def\bnodeheight{1.618*\bnodewidth}
\def\bcommitnodewidth{\bnodewidth}
\def\bcommitnodeheight{\bnodeheight}
\def\rnodewidth{(1/8)*\bnodewidth}
\def\rnodeheight{\rnodewidth}

% common styles
\tikzset{
    % small rbyd tag/alt styles
    rtag/.style={
        draw,
        ultra thick,
        minimum width=0.6*\rnodewidth,
        minimum height=0.6*\rnodeheight,
        inner sep=0,
        outer sep=0,
        fill=fglightblue,
    },
    ralt/.style={rtag, circle},
    raltb/.style={ralt},
    raltr/.style={ralt, draw=fgdarkred},
    ralty/.style={ralt, draw=fgdarkyellow},
    raltptr/.style={-{To[scale=0.5]}, ultra thick},
    raltbptr/.style={raltptr},
    raltrptr/.style={raltptr, draw=fgdarkred},
    raltyptr/.style={raltptr, draw=fgdarkyellow},
    % btree styles
    btree/.style={
        draw,
        ultra thick,
        fill=fglightblue,
        minimum width=\bnodewidth,
        minimum height=\bnodeheight,
        inner sep=0,
        outer sep=0,
    },
    bptr/.style={
        ->,
        %ultra thick,
        line width=2.4*1.6pt,
    },
    bdashedptr/.style={
        bptr,
        %loosely dashed,
        dash pattern=on 6pt off 6pt,
        shorten <= 0.125*\bnodewidth,
        shorten >= 0.125*\bnodewidth,
    },
    boutline/.style={
        draw,
        ultra thick,
    },
    bdashedoutline/.style={
        boutline,
        %loosely dashed,
        dash pattern=on 6pt off 6pt,
    },
    bbubble/.style={
        draw,
        ultra thick,
        %loosely dashed,
        dash pattern=on 6pt off 6pt,
        rounded corners=40pt,
        inner sep=0,
        outer sep=0,
    },
    bbubbleptr/.style={
        ->,
        %ultra thick,
        line width=2.4*1.6pt,
        %loosely dashed,
        dash pattern=on 6pt off 6pt,
        shorten <= 0.5*\bnodewidth,
        shorten >= 0.5*\bnodewidth,
    },
}


% no I'm not reimplementing rbyd trees in tikz
\newcommand{\rbydvi}[1]{
    \begin{scope}[
        x=\rnodewidth,
        y=\rnodeheight,
        node distance=1,
        on grid,
    ]
        % test node
        %\node at (#1.center) [rtag] {};

        % trunk 0
        \begin{scope}[on above layer]
        \node (#1_r0_0) [rtag, below left=2.5 and 3 of #1] {};
        \end{scope}

        % trunk 1
        \begin{scope}[on above layer]
        \node (#1_r1_0) [rtag, right=1 of #1_r0_0] {};
        \node (#1_r1_1) [raltb, above=1 of #1_r1_0] {};
        \end{scope}
        \draw[raltbptr] (#1_r1_1) -- (#1_r1_0);
        \draw[raltbptr] (#1_r1_1) -- (#1_r0_0);

        % trunk 2
        \begin{scope}[on above layer]
        \node (#1_r2_0) [rtag, right=2 of #1_r0_0] {};
        \node (#1_r2_1) [raltb, above=1 of #1_r2_0] {};
        \node (#1_r2_2) [raltr, above=2 of #1_r2_0] {};
        \end{scope}
        \draw[raltbptr] (#1_r2_1) -- (#1_r2_0);
        \draw[raltbptr] (#1_r2_1) -- (#1_r1_0);
        \draw[raltrptr] (#1_r2_2) -- (#1_r2_1);
        \draw[raltbptr] (#1_r2_2) -- (#1_r0_0);

        % trunk 3
        \begin{scope}[on above layer]
        \node (#1_r3_0) [rtag, right=3 of #1_r0_0] {};
        \node (#1_r3_1) [raltb, above=1 of #1_r3_0] {};
        \node (#1_r3_2) [raltr, above=2 of #1_r3_0] {};
        \node (#1_r3_3) [ralty, above=3 of #1_r3_0] {};
        \end{scope}
        \draw[raltbptr] (#1_r3_1) -- (#1_r3_0);
        \draw[raltbptr] (#1_r3_1) -- (#1_r2_0);
        \draw[raltrptr] (#1_r3_2) -- (#1_r3_1);
        \draw[raltbptr] (#1_r3_2) -- (#1_r1_0);
        \draw[raltyptr] (#1_r3_3) -- (#1_r3_2);
        \draw[raltbptr] (#1_r3_3) -- (#1_r0_0);

        % trunk 4
        \begin{scope}[on above layer]
        \node (#1_r4_0) [rtag, right=4 of #1_r0_0] {};
        \node (#1_r4_1) [raltb, above=1 of #1_r4_0] {};
        \node (#1_r4_2) [raltr, above=2 of #1_r4_0] {};
        \node (#1_r4_3) [raltb, above=3 of #1_r4_0] {};
        \end{scope}
        \draw[raltbptr] (#1_r4_1) -- (#1_r4_0);
        \draw[raltbptr] (#1_r4_1) -- (#1_r3_0);
        \draw[raltrptr] (#1_r4_2) -- (#1_r4_1);
        \draw[raltbptr] (#1_r4_2) -- (#1_r2_0);
        \draw[raltbptr] (#1_r4_3) -- (#1_r4_2);
        \draw[raltbptr] (#1_r4_3) -- (#1_r3_3);

        % trunk 5
        \begin{scope}[on above layer]
        \node (#1_r5_0) [rtag, right=5 of #1_r0_0] {};
        \node (#1_r5_1) [raltb, above=1 of #1_r5_0] {};
        \node (#1_r5_2) [raltr, above=2 of #1_r5_0] {};
        \node (#1_r5_3) [ralty, above=3 of #1_r5_0] {};
        \node (#1_r5_4) [raltb, above=4 of #1_r5_0] {};
        \end{scope}
        \draw[raltbptr] (#1_r5_1) -- (#1_r5_0);
        \draw[raltbptr] (#1_r5_1) -- (#1_r4_0);
        \draw[raltrptr] (#1_r5_2) -- (#1_r5_1);
        \draw[raltbptr] (#1_r5_2) -- (#1_r3_0);
        \draw[raltyptr] (#1_r5_3) -- (#1_r5_2);
        \draw[raltbptr] (#1_r5_3) -- (#1_r2_0);
        \draw[raltbptr] (#1_r5_4) -- (#1_r5_3);
        \draw[raltbptr] (#1_r5_4) -- (#1_r3_3);

        % trunk 6
        \begin{scope}[on above layer]
        \node (#1_r6_0) [rtag, right=6 of #1_r0_0] {};
        \node (#1_r6_1) [raltb, above=1 of #1_r6_0] {};
        \node (#1_r6_2) [raltr, above=2 of #1_r6_0] {};
        \node (#1_r6_3) [raltb, above=3 of #1_r6_0] {};
        \node (#1_r6_4) [raltr, above=4 of #1_r6_0, alias=#1_rroot] {};
        \end{scope}
        \draw[raltbptr] (#1_r6_1) -- (#1_r6_0);
        \draw[raltbptr] (#1_r6_1) -- (#1_r5_0);
        \draw[raltrptr] (#1_r6_2) -- (#1_r6_1);
        \draw[raltbptr] (#1_r6_2) -- (#1_r4_0);
        \draw[raltbptr] (#1_r6_3) -- (#1_r6_2);
        \draw[raltbptr] (#1_r6_3) -- (#1_r5_3);
        \draw[raltrptr] (#1_r6_4) -- (#1_r6_3);
        \draw[raltbptr] (#1_r6_4) -- (#1_r3_3);
    \end{scope}
}

\newcommand{\rbydreplacevii}[2]{
    \begin{scope}[
        x=\rnodewidth,
        y=\rnodeheight,
        node distance=1,
        on grid,
    ]
        % test node
        %\node at (#2.center) [rtag] {};

        % trunk 7
        \begin{scope}[on above layer]
        \node (#2_r7_0) [rtag, below=2.5 of #2] {};
        \node (#2_r7_1) [raltb, above=1 of #2_r7_0] {};
        \node (#2_r7_2) [raltr, above=2 of #2_r7_0] {};
        \node (#2_r7_3) [raltb, above=3 of #2_r7_0] {};
        \node (#2_r7_4) [raltr, above=4 of #2_r7_0, alias=#2_rroot] {};
        \end{scope}
        \draw[raltbptr] (#2_r7_1) -- (#2_r7_0);
        \draw[raltbptr] (#2_r7_1) -- (#1_r5_0);
        \draw[raltrptr] (#2_r7_2) -- (#2_r7_1);
        \draw[raltbptr] (#2_r7_2) -- (#1_r4_0);
        \draw[raltbptr] (#2_r7_3) -- (#2_r7_2);
        \draw[raltbptr] (#2_r7_3) -- (#1_r5_3);
        \draw[raltrptr] (#2_r7_4) -- (#2_r7_3);
        \draw[raltbptr] (#2_r7_4) -- (#1_r3_3);
    \end{scope}
}

\newcommand{\rbydappendviii}[2]{
    \begin{scope}[
        x=\rnodewidth,
        y=\rnodeheight,
        node distance=1,
        on grid,
    ]
        % test node
        %\node at (#2.center) [rtag] {};

        % trunk 7
        \begin{scope}[on above layer]
        \node (#2_r7_0) [rtag, below=2.5 of #2] {};
        \node (#2_r7_1) [raltb, above=1 of #2_r7_0] {};
        \node (#2_r7_2) [raltr, above=2 of #2_r7_0] {};
        \node (#2_r7_3) [ralty, above=3 of #2_r7_0] {};
        \node (#2_r7_4) [raltb, above=4 of #2_r7_0] {};
        \node (#2_r7_5) [raltr, above=5 of #2_r7_0, alias=#2_rroot] {};
        \end{scope}
        \draw[raltbptr] (#2_r7_1) -- (#2_r7_0);
        \draw[raltbptr] (#2_r7_1) -- (#1_r6_0);
        \draw[raltrptr] (#2_r7_2) -- (#2_r7_1);
        \draw[raltbptr] (#2_r7_2) -- (#1_r5_0);
        \draw[raltyptr] (#2_r7_3) -- (#2_r7_2);
        \draw[raltbptr] (#2_r7_3) -- (#1_r4_0);
        \draw[raltbptr] (#2_r7_4) -- (#2_r7_3);
        \draw[raltbptr] (#2_r7_4) -- (#1_r5_3);
        \draw[raltrptr] (#2_r7_5) -- (#2_r7_4);
        \draw[raltbptr] (#2_r7_5) -- (#1_r3_3);
    \end{scope}
}

\newcommand{\rbydreplaceix}[2]{
    \begin{scope}[
        x=\rnodewidth,
        y=\rnodeheight,
        node distance=1,
        on grid,
    ]
        % test node
        %\node at (#2.center) [rtag] {};

        % trunk 7
        \begin{scope}[on above layer]
        \node (#2_r7_0) [rtag, below left=2.5 and 0.5 of #2] {};
        \node (#2_r7_1) [raltb, above=1 of #2_r7_0] {};
        \node (#2_r7_2) [raltr, above=2 of #2_r7_0] {};
        \node (#2_r7_3) [raltb, above=3 of #2_r7_0] {};
        \node (#2_r7_4) [raltr, above=4 of #2_r7_0] {};
        \end{scope}
        \draw[raltbptr] (#2_r7_1) -- (#2_r7_0);
        \draw[raltbptr] (#2_r7_1) -- (#1_r5_0);
        \draw[raltrptr] (#2_r7_2) -- (#2_r7_1);
        \draw[raltbptr] (#2_r7_2) -- (#1_r4_0);
        \draw[raltbptr] (#2_r7_3) -- (#2_r7_2);
        \draw[raltbptr] (#2_r7_3) -- (#1_r5_3);
        \draw[raltrptr] (#2_r7_4) -- (#2_r7_3);
        \draw[raltbptr] (#2_r7_4) -- (#1_r3_3);

        % trunk 8
        \begin{scope}[on above layer]
        \node (#2_r8_0) [rtag, right=1 of #2_r7_0] {};
        \node (#2_r8_1) [raltb, above=1 of #2_r8_0] {};
        \node (#2_r8_2) [raltr, above=2 of #2_r8_0] {};
        \node (#2_r8_3) [ralty, above=3 of #2_r8_0] {};
        \node (#2_r8_4) [raltb, above=4 of #2_r8_0] {};
        \node (#2_r8_5) [raltr, above=5 of #2_r8_0, alias=#2_rroot] {};
        \end{scope}
        \draw[raltbptr] (#2_r8_1) -- (#2_r8_0);
        \draw[raltbptr] (#2_r8_1) -- (#2_r7_0);
        \draw[raltrptr] (#2_r8_2) -- (#2_r8_1);
        \draw[raltbptr] (#2_r8_2) -- (#1_r5_0);
        \draw[raltyptr] (#2_r8_3) -- (#2_r8_2);
        \draw[raltbptr] (#2_r8_3) -- (#1_r4_0);
        \draw[raltbptr] (#2_r8_4) -- (#2_r8_3);
        \draw[raltbptr] (#2_r8_4) -- (#1_r5_3);
        \draw[raltrptr] (#2_r8_5) -- (#2_r8_4);
        \draw[raltbptr] (#2_r8_5) -- (#1_r3_3);
    \end{scope}
}

\newcommand{\rbydrebalancevi}[1]{
    \begin{scope}[
        x=\rnodewidth,
        y=\rnodeheight,
        node distance=1,
        on grid,
    ]
        % test node
        %\node at (#1.center) [rtag] {};

        % trunk 0
        \begin{scope}[on above layer]
        \node (#1_r0_0) [rtag, below left=2.5 and 3 of #1] {};
        \end{scope}

        % trunk 1
        \begin{scope}[on above layer]
        \node (#1_r1_0) [rtag, right=1 of #1_r0_0] {};
        \node (#1_r1_1) [raltb, above=1 of #1_r1_0] {};
        \end{scope}
        \draw[raltbptr] (#1_r1_1) -- (#1_r1_0);
        \draw[raltbptr] (#1_r1_1) -- (#1_r0_0);

        % trunk 2
        \begin{scope}[on above layer]
        \node (#1_r2_0) [rtag, right=2 of #1_r0_0] {};
        \end{scope}

        % trunk 3
        \begin{scope}[on above layer]
        \node (#1_r3_0) [rtag, right=3 of #1_r0_0] {};
        \node (#1_r3_1) [raltb, above=1 of #1_r3_0] {};
        \node (#1_r3_2) [raltb, above=2 of #1_r3_0] {};
        \end{scope}
        \draw[raltbptr] (#1_r3_1) -- (#1_r3_0);
        \draw[raltbptr] (#1_r3_1) -- (#1_r2_0);
        \draw[raltbptr] (#1_r3_2) -- (#1_r3_1);
        \draw[raltbptr] (#1_r3_2) -- (#1_r1_1);

        % trunk 4
        \begin{scope}[on above layer]
        \node (#1_r4_0) [rtag, right=4 of #1_r0_0] {};
        \end{scope}

        % trunk 5
        \begin{scope}[on above layer]
        \node (#1_r5_0) [rtag, right=5 of #1_r0_0] {};
        \node (#1_r5_1) [raltb, above=1 of #1_r5_0] {};
        \end{scope}
        \draw[raltbptr] (#1_r5_1) -- (#1_r5_0);
        \draw[raltbptr] (#1_r5_1) -- (#1_r4_0);

        % trunk 6
        \begin{scope}[on above layer]
        \node (#1_r6_0) [rtag, right=6 of #1_r0_0] {};
        \node (#1_r6_1) [raltb, above=1 of #1_r6_0] {};
        \node (#1_r6_2) [raltb, above=2 of #1_r6_0] {};
        \node (#1_r6_3) [raltb, above=3 of #1_r6_0, alias=#1_rroot] {};
        \end{scope}
        \draw[raltbptr] (#1_r6_1) -- (#1_r6_0);
        \draw[raltbptr] (#1_r6_2) -- (#1_r6_1);
        \draw[raltbptr] (#1_r6_2) -- (#1_r5_1);
        \draw[raltbptr] (#1_r6_3) -- (#1_r6_2);
        \draw[raltbptr] (#1_r6_3) -- (#1_r3_2);
    \end{scope}
}

\newcommand{\rbydrebalanceiv}[1]{
    \begin{scope}[
        x=\rnodewidth,
        y=\rnodeheight,
        node distance=1,
        on grid,
    ]
        % test node
        %\node at (#1.center) [rtag] {};

        % trunk 0
        \begin{scope}[on above layer]
        \node (#1_r0_0) [rtag, below left=2.5 and 1.5 of #1] {};
        \end{scope}

        % trunk 1
        \begin{scope}[on above layer]
        \node (#1_r1_0) [rtag, right=1 of #1_r0_0] {};
        \node (#1_r1_1) [raltb, above=1 of #1_r1_0] {};
        \end{scope}
        \draw[raltbptr] (#1_r1_1) -- (#1_r1_0);
        \draw[raltbptr] (#1_r1_1) -- (#1_r0_0);

        % trunk 2
        \begin{scope}[on above layer]
        \node (#1_r2_0) [rtag, right=2 of #1_r0_0] {};
        \end{scope}

        % trunk 3
        \begin{scope}[on above layer]
        \node (#1_r3_0) [rtag, right=3 of #1_r0_0] {};
        \node (#1_r3_1) [raltb, above=1 of #1_r3_0] {};
        \node (#1_r3_2) [raltb, above=2 of #1_r3_0, alias=#1_rroot] {};
        \end{scope}
        \draw[raltbptr] (#1_r3_1) -- (#1_r3_0);
        \draw[raltbptr] (#1_r3_1) -- (#1_r2_0);
        \draw[raltbptr] (#1_r3_2) -- (#1_r3_1);
        \draw[raltbptr] (#1_r3_2) -- (#1_r1_1);
    \end{scope}
}

\newcommand{\rbydrebalanceiii}[1]{
    \begin{scope}[
        x=\rnodewidth,
        y=\rnodeheight,
        node distance=1,
        on grid,
    ]
        % test node
        %\node at (#1.center) [rtag] {};

        % trunk 0
        \begin{scope}[on above layer]
        \node (#1_r0_0) [rtag, below left=2.5 and 1 of #1] {};
        \end{scope}

        % trunk 1
        \begin{scope}[on above layer]
        \node (#1_r1_0) [rtag, right=1 of #1_r0_0] {};
        \node (#1_r1_1) [raltb, above=1 of #1_r1_0] {};
        \end{scope}
        \draw[raltbptr] (#1_r1_1) -- (#1_r1_0);
        \draw[raltbptr] (#1_r1_1) -- (#1_r0_0);

        % trunk 2
        \begin{scope}[on above layer]
        \node (#1_r2_0) [rtag, right=2 of #1_r0_0] {};
        \node (#1_r2_1) [raltb, above=1 of #1_r2_0] {};
        \node (#1_r2_2) [raltb, above=2 of #1_r2_0, alias=#1_rroot] {};
        \end{scope}
        \draw[raltbptr] (#1_r2_1) -- (#1_r2_0);
        \draw[raltbptr] (#1_r2_2) -- (#1_r2_1);
        \draw[raltbptr] (#1_r2_2) -- (#1_r1_1);
    \end{scope}
}


% B-tree commit example
\newcommand{\rbydbtreecommit}[3]{
    \begin{scope}[
        x=\bcommitnodewidth,
        y=\bcommitnodeheight,
        node distance=1,
        on grid,
    ]
        % containing bubble
        \node (#1) [
            bbubble,
            #2 of #3,
            minimum width=3.25*\bcommitnodewidth,
            minimum height=3.125*\bcommitnodeheight] {};

        % title
        \node at (#1.north) [
            anchor=north,
            %minimum height=0.25*\bcommitnodeheight,
            minimum height=0.33*\bcommitnodeheight,
        ] {\Large
            Incremental commit
        };

        % B-tree nodes
        \node (bc1_0) [
            btree,
            above left=2.25 and 0.75 of #1.south,
            anchor=center,
            draw=none
        ] {};
        \rbydvi{bc1_0};

        \node (bc0_0) [
            btree,
            below right=1.5 and 0.5 of bc1_0,
            draw=none,
        ] {};
        \rbydvi{bc0_0};

        \draw[bptr]
            (bc1_0_r6_0)
            .. controls +(0,-1) and +(0,+1) ..
            (bc0_0_rroot);

        % B-tree commit
        \node (bc1_1) [
            btree,
            right=1.25 of bc1_0,
            minimum width=0.5*\bcommitnodewidth,
            draw=none,
        ] {};
        \scoped[on behind layer] \draw[fill=fglightblue!50!white]
            (bc1_0.south west) rectangle (bc1_1.north east);
        \draw[boutline]
            (bc1_0.north east)
            -- (bc1_0.north west)
            -- (bc1_0.south west)
            -- (bc1_0.south east)
            (bc1_1.south west)
            -- (bc1_1.south east)
            -- (bc1_1.north east)
            -- (bc1_1.north west);
        \draw[bdashedoutline]
            (bc1_0.south east)
            -- (bc1_1.south west)
            (bc1_1.north west)
            -- (bc1_0.north east);
        \rbydreplacevii{bc1_0}{bc1_1};

        \node (bc0_1) [
            btree,
            right=1.25 of bc0_0,
            minimum width=0.5*\bcommitnodewidth,
            draw=none,
        ] {};
        \scoped[on behind layer] \draw[fill=fglightblue!50!white]
            (bc0_0.south west) rectangle (bc0_1.north east);
        \draw[boutline]
            (bc0_0.north east)
            -- (bc0_0.north west)
            -- (bc0_0.south west)
            -- (bc0_0.south east)
            (bc0_1.south west)
            -- (bc0_1.south east)
            -- (bc0_1.north east)
            -- (bc0_1.north west);
        \draw[bdashedoutline]
            (bc0_0.south east)
            -- (bc0_1.south west)
            (bc0_1.north west)
            -- (bc0_0.north east);
        \rbydappendviii{bc0_0}{bc0_1};

        \draw[bptr]
            (bc1_1_r7_0)
            .. controls +(0,-1) and +(0,+1) ..
            (bc0_1_rroot);

        % op arrow
        \draw[bbubbleptr] (#3) -- (#1)
            node[pos=0.5, anchor=south] {\Large
                $\text{uncompacted} \le b$
            };
    \end{scope}
}

% B-tree compact example
\newcommand{\rbydbtreecompact}[3]{
    \begin{scope}[
        x=\bcommitnodewidth,
        y=\bcommitnodeheight,
        node distance=1,
        on grid,
    ]
        % containing bubble
        \node (#1) [
            bbubble,
            #2 of #3,
            minimum width=3.5*\bcommitnodewidth,
            minimum height=3.125*\bcommitnodeheight] {};

        % title
        \node at (#1.north) [
            anchor=north,
            %minimum height=0.25*\bcommitnodeheight,
            minimum height=0.33*\bcommitnodeheight,
        ] {\Large
            Compact 
        };

        % B-tree nodes
        \node (bc1_0) [
            btree,
            above left=2.25 and 0.625 of #1.south,
            anchor=center,
            draw=none
        ] {};
        \rbydvi{bc1_0};

        \node (bc0_0) [
            btree,
            below right=1.5 and -0.25 of bc1_0,
        ] {};
        \rbydvi{bc0_0};

        \draw[bptr]
            (bc1_0_r6_0)
            .. controls +(0,-1) and +(0,+1) ..
            (bc0_0_rroot);

        % B-tree commit
        \node (bc1_1) [
            btree,
            right=1.25 of bc1_0,
            minimum width=0.5*\bcommitnodewidth,
            draw=none,
        ] {};
        \scoped[on behind layer] \draw[fill=fglightblue!50!white]
            (bc1_0.south west) rectangle (bc1_1.north east);
        \draw[boutline]
            (bc1_0.north east)
            -- (bc1_0.north west)
            -- (bc1_0.south west)
            -- (bc1_0.south east)
            (bc1_1.south west)
            -- (bc1_1.south east)
            -- (bc1_1.north east)
            -- (bc1_1.north west);
        \draw[bdashedoutline]
            (bc1_0.south east)
            -- (bc1_1.south west)
            (bc1_1.north west)
            -- (bc1_0.north east);
        \rbydreplacevii{bc1_0}{bc1_1};

        % B-tree compact
        \node (bc0_1) [
            btree,
            right=1.75 of bc0_0,
        ] {};
        \rbydrebalancevi{bc0_1};

        \draw[bptr]
            (bc1_1_r7_0)
            .. controls +(0,-1) and +(0,+1) ..
            (bc0_1_rroot);

        \draw[bdashedptr] (bc0_0) -- (bc0_1);

        % op arrow
        \draw[bbubbleptr] (#3) -- (#1)
            node[pos=0.5, anchor=south] {\Large
                %$\text{compacted} \le b/2$
                $\text{compacted} \le \frac{b}{2}$
            };
    \end{scope}
}

% B-tree split example
\newcommand{\rbydbtreesplit}[3]{
    \begin{scope}[
        x=\bcommitnodewidth,
        y=\bcommitnodeheight,
        node distance=1,
        on grid,
    ]
        % containing bubble
        \node (#1) [
            bbubble,
            #2 of #3,
            minimum width=4.75*\bcommitnodewidth,
            minimum height=3.125*\bcommitnodeheight] {};

        % title
        \node at (#1.north) [
            anchor=north,
            %minimum height=0.25*\bcommitnodeheight,
            minimum height=0.33*\bcommitnodeheight,
        ] {\Large
            Split 
        };

        % B-tree nodes
        \node (bc1_0) [
            btree,
            above left=2.25 and 0.625 of #1.south,
            anchor=center,
            draw=none
        ] {};
        \rbydvi{bc1_0};

        \node (bc0_0) [
            btree,
            below right=1.5 and -0.875 of bc1_0,
        ] {};
        \rbydvi{bc0_0};

        \draw[bptr]
            (bc1_0_r6_0)
            .. controls +(0,-1) and +(0,+1) ..
            (bc0_0_rroot);

        % B-tree commit
        \node (bc1_1) [
            btree,
            right=1.25 of bc1_0,
            minimum width=0.5*\bcommitnodewidth,
            draw=none,
        ] {};
        \scoped[on behind layer] \draw[fill=fglightblue!50!white]
            (bc1_0.south west) rectangle (bc1_1.north east);
        \draw[boutline]
            (bc1_0.north east)
            -- (bc1_0.north west)
            -- (bc1_0.south west)
            -- (bc1_0.south east)
            (bc1_1.south west)
            -- (bc1_1.south east)
            -- (bc1_1.north east)
            -- (bc1_1.north west);
        \draw[bdashedoutline]
            (bc1_0.south east)
            -- (bc1_1.south west)
            (bc1_1.north west)
            -- (bc1_0.north east);
        \rbydreplaceix{bc1_0}{bc1_1};

        % B-tree split
        \node (bc0_1) [
            btree,
            right=1.75 of bc0_0,
        ] {};
        \rbydrebalanceiv{bc0_1};
        \node (bc0_2) [
            btree,
            right=1.25 of bc0_1,
        ] {};
        \rbydrebalanceiii{bc0_2};

        \draw[bptr]
            (bc1_1_r7_0)
            .. controls +(0,-1) and +(0,+1) ..
            (bc0_1_rroot);
        \draw[bptr]
            (bc1_1_r8_0)
            .. controls +(0,-1) and +(0,+1.1) ..
            (bc0_2_rroot);

        \draw[bdashedptr] (bc0_0) -- (bc0_1);

        % op arrow
        \draw[bbubbleptr] (#3) -- (#1)
            node[pos=0.5, anchor=south] {\Large
                %$\text{compacted} \gt b/2$
                $\text{compacted} \gt \frac{b}{2}$
            };
    \end{scope}
}



% main B-tree
\newcommand{\rbydbtree}{
    \begin{scope}[
        x=\bnodewidth,
        y=\bnodeheight,
        node distance=1,
        on grid,
    ]
        % B-tree nodes
        \node (b0_0) [btree] {};
        \rbydvi{b0_0};
        \node (b0_1) [btree, right=1*1.5 of b0_0] {};
        \rbydvi{b0_1};
        \node (b0_2) [btree, right=2*1.5 of b0_0] {};
        \rbydvi{b0_2};
        \node (b0_3) [btree, right=3*1.5 of b0_0] {};
        \rbydvi{b0_3};
        \node (b0_4) [btree, right=4*1.5 of b0_0] {};
        \rbydvi{b0_4};
        \node (b0_5) [btree, right=5*1.5 of b0_0] {};
        \rbydvi{b0_5};
        \node (b0_6) [btree, right=6*1.5 of b0_0] {};
        \rbydvi{b0_6};

        \node (b1_0) [btree, above right=1.5 and 4.5 of b0_0, alias=broot] {};
        \rbydvi{b1_0};
        \draw[bptr]
            (b1_0_r0_0)
            .. controls +(0,-1) and +(0,+1.3) ..
            (b0_0_rroot);
        \draw[bptr]
            (b1_0_r1_0)
            .. controls +(0,-1) and +(0,+1.2) ..
            (b0_1_rroot);
        \draw[bptr]
            (b1_0_r2_0)
            .. controls +(0,-1) and +(0,+1.1) ..
            (b0_2_rroot);
        \draw[bptr]
            (b1_0_r3_0)
            .. controls +(0,-1) and +(0,+1) ..
            (b0_3_rroot);
        \draw[bptr]
            (b1_0_r4_0)
            .. controls +(0,-1) and +(0,+1.1) ..
            (b0_4_rroot);
        \draw[bptr]
            (b1_0_r5_0)
            .. controls +(0,-1) and +(0,+1.2) ..
            (b0_5_rroot);
        \draw[bptr]
            (b1_0_r6_0)
            .. controls +(0,-1) and +(0,+1.3) ..
            (b0_6_rroot);

        % B-tree commit/compact/split examples
        \rbydbtreecommit{bcommit}{right=6}{broot};
        \rbydbtreecompact{bcompact}{above right=1 and 3.75}{bcommit};
        \rbydbtreesplit{bsplit}{below right=0.5 and 7}{bcommit};

    \end{scope}
}


% put big B-tree mess into our poster
\block[bodyinnersep=0pt]{}{
    \begin{tikzfigure}
    \flushleft
    \begin{tikzpicture}
        \rbydbtree
    \end{tikzpicture}
    \end{tikzfigure}
}



% B-tree details
\begin{columns}
\column{0.33}
\block{Rbyd B-trees}{\Large
}
\column{0.33}
\block{Deferred splits/merges}{\Large
}
\column{0.33}
\block{Rebalancing compactions}{\Large
}
\end{columns}


% Results
\block{Results}{\Large
}




% \block{}
% {
% \blindtext
% }
% 
% \begin{columns}
% \column{0.4}
% \block{More text}{Text and more text}
% 
% \column{0.6}
% \block{Something else}{Here, \blindtext \vspace{4cm}}
% \note[%targetoffsetx=-9cm, targetoffsety=-5.5cm%, width=0.5\linewidth
% ]{This is a note...}
% \end{columns}
% 
% \begin{columns}
% \column{0.5}
% \block{A figure}
% {
% \begin{center}
%     %\includegraphics[width=0.2\textheight]{logo.png}
%     \input{littlefs-ico}
% \end{center}
% }
% 
% \block{Outer Block}{
% \blindtext
% \vspace{24pt}
% \innerblock[]{Inner Block}{\blindtext}
% }
% 
% \column{0.5}
% \block{Description of the figure}{\Large\blindtext[2]}
% \end{columns}

\end{document}

